% 原书第 120—122 页；6.8。
\section{因果模型的等价性与可否证性}
\label{sec:ch06-equivalence}

到目前为止，结构因果模型一直是数学对象。为将其与现实联系起来，我们把它视为数据生成过程的模型。不过，这可能是一类复杂的模型。现在，我们不仅仅对一个联合分布建模——例如用 Poisson 过程对某个物理过程建模——还可以同时对系统处于观测状态以及受到扰动时的情况建模。前面已经看到，甚至可以将结构因果模型视为反事实陈述的模型。

更形式化地说，考虑随机变量向量 $\mathbf X=(X_1,\ldots,X_d)$。$\mathbf X$ 的\textbf{概率模型（probabilistic model）}预测观测分布 $P_{\mathbf X}$。如果模型还能预测将某些变量 $X_j$ 设为独立变量 $\widetilde N_j$ 后的干预分布，就称其为\textbf{干预模型（interventional model）}。最后，\textbf{反事实模型（counterfactual model）}还能够预测反事实陈述的结果。例如，传统机器学习方法建立概率模型；因果图模型（定义 6.32）可用作干预模型；结构因果模型则可用作反事实模型。若两个模型的对应预测相同，就称其等价；类似构造见 Bongers et al.［2016］。

\Needspace{4cm}
\begin{definition}[因果模型的等价性]
若两个模型导出相同的观测分布，就称为\textbf{概率等价（probabilistically equivalent）}；若导出相同的观测分布与干预分布，就称为\textbf{干预等价（interventionally equivalent）}；若导出相同的观测分布、干预分布与反事实分布，就称为\textbf{反事实等价（counterfactually equivalent）}。
\end{definition}

显然，例如干预等价的概念，只适用于干预模型与反事实模型。命题 7.1 意味着，对每个概率模型，都存在一个与之观测等价的结构因果模型。

若 $\mathbf X$ 的密度严格为正，命题 6.48 表明，可以只考察单节点干预：将某个变量 $X_j$ 设为分布具有完整支撑集的变量 $\widetilde N_j$。如果两个模型在这类干预上相同，那么在所有其他干预上也相同。这样做的理由是，单节点干预对应标准形式的随机化实验。

对于给定的数据生成过程，若概率模型或干预模型给出的相应分布，与从该过程观测到的数据不一致，就可以\textbf{否证（falsify）}模型。也就是说，即使一个干预模型正确预测了观测分布，但无法预测随机化实验中会发生什么，它仍被视为已被否证。这一概念包含一项假设：对于随机化实验应当是什么样，存在共同认可的理解。如果不清楚现实中如何对相关变量进行随机化或干预，就应谨慎地写出结构因果模型。可否证性还需要统计显著性的概念，这里不作讨论。

我们没有将反事实模型纳入上述讨论，因为它们通常难以否证。可以根据它们对观测分布与干预分布的推论来否证它们，见 Shpitser and Pearl［2008a］及其参考文献。在某些特定实验设计下，还可以构造可否证的反事实陈述（见例 3.4）。但例 6.19 也说明，两个结构因果模型可以导出相同的观测分布和干预分布，却导出不同的反事实陈述。

上面将干预限制为“把单个变量设为噪声变量”这一子类，是出于实践考虑。检验模型是否有效，需要比较随机化实验的结果与模型预测。对于更复杂的干预，现实中的对应实验似乎更难实施。下面的命题说明，这种限制并不损失一般性：只要因果模型在所有单节点干预上相同，它们就是干预等价的。证明见附录 C.7。

\begin{proposition}[干预等价性]
假设两个结构因果模型或因果图模型 $\mathfrak C_1$、$\mathfrak C_2$ 诱导出严格为正的连续条件密度 $p(x_j\mid x_{\pa(j)})$，其中 $\pa(j):=\PA_{X_j}$，并满足因果极小性。进一步假设：将任意变量 $X_j$ 设为具有完整支撑集的变量 $\widetilde N_j$ 时，两个模型导出的干预分布相同：
\[
 P_{\mathbf X}^{\mathfrak C_1;\dox{X_j:=\widetilde N_j}}
 =P_{\mathbf X}^{\mathfrak C_2;\dox{X_j:=\widetilde N_j}},
 \qquad\forall j,\ \forall\widetilde N_j\text{ 具有完整支撑集}.
\]
则 $\mathfrak C_1$ 与 $\mathfrak C_2$ 干预等价，即对任意可能的干预都相同，包括原子干预，以及改变父节点集但不产生环的干预。
\end{proposition}

如果密度并非严格为正，这一结论就不一定成立，此时可能需要考虑同时干预多个节点，例如双基因敲除实验；见习题 6.59。

此外，现在可以进一步论证结构因果模型的结构极小性概念，见注 6.6。前面指出，如果结构赋值中的函数不依赖某个输入，可以选择更稀疏的表示。下面的命题将这些表示在哪种意义下等价加以形式化。

\Needspace{3cm}
\begin{proposition}[反事实等价性]
考虑两个结构因果模型 $\mathfrak C$ 和 $\mathfrak C^*$。它们具有相同的噪声分布 $P_{\mathbf N}$，仅第 $k$ 个结构赋值不同：
\begin{equation}
 f_k(\pa_k,n_k)=f_k^*(\pa_k^*,n_k),\qquad
 \forall\pa_k,\ \forall n_k\text{ 满足 }p(n_k)>0,
 \tag{6.24}
\end{equation}
其中 $\PA_k^*\subsetneq\PA_k$。则两个模型反事实等价。
\end{proposition}
证明见附录 C.8。
\FloatBarrier
